%%%%%%%%%%%%%%%%%%%%%%%%%%%%%%%%%%%%%%%%%%%%%%%%%%
% 2.EINSTELLUNGEN
%%%%%%%%%%%%%%%%%%%%%%%%%%%%%%%%%%%%%%%%%%%%%%%%%%

% Indexgenerierung (Prozessoranweisung)
\makeindex

% setzt den Pfad für Graphiken
\graphicspath{{figures/}{content/img/}{deployment/}{deployment/screens/}{img/}}
% zeigt Nummerierungen bei subsubsections an
\setcounter{secnumdepth}{4}
% Zähler der die maximale Nummer von "Floatobjekten" am Beginn einer Seite angibt
\setcounter{topnumber}{2}
% Redefinition: gibt an, wieviel Platz Floatobjekte (Tabellen, Bilder)
% am Beginn einer Seite einnehmen dürfen. (80 Prozent)
\def\topfraction{.8}
% Zähler der die maximale Nummer von "Floatobjekten" am Ende einer Seite angibt
\setcounter{bottomnumber}{2}
% Redefinition: gibt an, wieviel Platz Floatobjekte (Tabellen, Bilder)
% am Beginn einer Seite einnehmen dürfen. (80 Prozent)
\def\bottomfraction{.5}
% maximale Anzahl der Floatobjekte auf einer Seite
\setcounter{totalnumber}{8}
% Redefinition: minimaler Prozentsatz einer Seite der Text auf dem Text stehen muss
\def\textfraction{.2}
% Redefinition: mimimaler Prozentsatz von Floatobjekten die auf einer Floatseite(Seite,
% die nur Tabellen, Bilder enthält) sein müssen
\def\floatpagefraction{.6}
% Einzug des Absatzes auf 0 pt stellen
\setlength{\parindent}{0pt}
% vertikaler Zwischenraum zwischen zwei Absätzen
\setlength{\parskip}{1ex plus 0.5ex minus 0.2ex}
% Teil des "caption" Packages: extra 20pkt links und rechts von einer Beschriftung
\setlength{\captionmargin}{20pt}
%Teil des "float" Packages
\floatstyle{plain}
% Name eines neuen "schwebenden" Objekts
\floatname{example}{Example}

\newfloat{example}{hbtp}{loe}[chapter]
\floatplacement{figure}{H}
\floatplacement{table}{H}
% Weniger Abstand um schwebende Objekte (reduziert grosse Luecken)
\setlength{\intextsep}{8pt plus 2pt minus 2pt}
\setlength{\textfloatsep}{10pt plus 2pt minus 2pt}

% transformiert "\dollar" zum Dollarzeichen
\newcommand{\dollar}{\char36}

% Bild mit abgerundeten Ecken (fuer App-Screenshots etc.)
\newcommand{\roundedimg}[2][]{%
  \begin{tikzpicture}[baseline=(img.base)]
    \node[rounded corners=4pt, inner sep=0, clip] (img) {\includegraphics[#1]{#2}};
  \end{tikzpicture}%
}	

% Skript fur die Abkürzungen
% neue Umgebung mit einem Argument
\newenvironment{bfscript}[1] {
 % Liste
 \begin{list}
 % keine Eintragsmarkierung
 {}
 {\settowidth{\labelwidth}{\bf #1}
  % Abstand vom Eintrag zum linken Rand auf die Breite des Labels setzen (0, da kein Label)
  \setlength{\leftmargin}{\labelwidth}
  % das Ganze um die Abstand des Label zum Text (auch 0) danach erhöhen
  \addtolength{\leftmargin}{\labelsep}
  % Abstand der Absätze innerhalb eines Eintrages
  \parsep{} 0.5ex plus 0.2ex minus 0.2ex
  % Abstand der einzelnen Einträge
  \itemsep{} 0.3ex
  % festlegen des Labels (fett, wenn nötig bis zum Text auffüllen)
  \renewcommand{\makelabel}[1]{\bf ##1\hfill}}}
 {\end{list}
}

% ============================================================
% Listings: TypeScript und Skyline-Stile (global fuer alle Kapitel)
% ============================================================
\lstdefinelanguage{TypeScript}{
  keywords={const,let,var,function,return,async,await,if,else,throw,
            new,class,import,from,export,default,try,catch,true,false,
            null,undefined,void,string,number,boolean,Promise},
  sensitive=true,
  comment=[l]{//},
  morecomment=[s]{/*}{*/},
  morestring=[b]',
  morestring=[b]"
}

\lstdefinestyle{skyline-sql}{
  language=SQL,
  basicstyle=\footnotesize\ttfamily,
  keywordstyle=\bfseries,
  commentstyle=\color{gray}\itshape,
  stringstyle=\color{mydarkelectricblue},
  numbers=left,
  numberstyle=\tiny\color{gray},
  numbersep=8pt,
  frame=tb,
  rulecolor=\color{mydarkelectricblue},
  breaklines=true,
  showstringspaces=false,
  captionpos=b,
  morekeywords={UUID,BIGINT,TIMESTAMPTZ,BOOLEAN,REFERENCES,CASCADE,
                POLICY,USING,CHECK,ENABLE,ROW,LEVEL,SECURITY,
                gen_random_uuid,auth,uid,CREATE,TABLE,IF,NOT,EXISTS,
                PRIMARY,KEY,DEFAULT,ON,DELETE,SET,NULL,INDEX,
                DROP,ALTER}
}

\lstdefinestyle{skyline-ts}{
  language=TypeScript,
  basicstyle=\footnotesize\ttfamily,
  keywordstyle=\bfseries,
  commentstyle=\color{gray}\itshape,
  stringstyle=\color{mydarkelectricblue},
  numbers=left,
  numberstyle=\tiny\color{gray},
  numbersep=8pt,
  frame=tb,
  rulecolor=\color{mydarkelectricblue},
  breaklines=true,
  showstringspaces=false,
  captionpos=b
}

% ============================================================
% Listings auf einer Seite halten (kein Umbruch mitten im Code)
% Bei weniger als 18cm freiem Platz -> Seitenumbruch vor dem Listing
% ============================================================
\BeforeBeginEnvironment{lstlisting}{\Needspace{18cm}}